% =====================================================================
\chapter{Υλοποίηση}
\label{ch:implementation}
% =====================================================================

\section{Τεχνολογικές επιλογές}
\label{sec:impl-tech}

Το \Probolos{} υλοποιήθηκε σε \en{Python}~3.10+ για τρεις λόγους. Πρώτον, η
\en{Python} είναι γλώσσα με ασφάλεια μνήμης: η ανάλυση μη έμπιστων δυαδικών
δεδομένων (περιγραφείς, τομείς δίσκου) δεν μπορεί να οδηγήσει σε υπερχείλιση
ενδιαμέσου. Δεύτερον, η βιβλιοθήκη προτύπων καλύπτει όλες τις απαιτούμενες
κλήσεις συστήματος (\code{socket.sendmsg} με \code{SCM\_RIGHTS},
\code{fcntl.ioctl}, \code{os.setgroups}, \code{prctl} μέσω \code{ctypes}).
Τρίτον, η μόνη εξωτερική εξάρτηση είναι το \en{pyudev} (απαίτηση~\ref{req:portable}),
ενώ το \en{PyYAML} είναι προαιρετικό για αρχεία κανόνων.

\section{Οργάνωση του κώδικα}
\label{sec:impl-modules}

Ο κώδικας παραγωγής αποτελείται από 35 αρχεία-μονάδες (\en{modules}) και περίπου 17\,700
γραμμές. Ο Πίνακας~\ref{tab:modules} συνοψίζει τις κυριότερες, ομαδοποιημένες
ανά ρόλο.

\begin{longtable}{@{}p{3.6cm}rp{8.6cm}@{}}
\caption{Οι κύριες μονάδες του \Probolos{} (γραμμές κώδικα στην έκδοση 1.0.0b2).}
\label{tab:modules}\\
\toprule
\textbf{Μονάδα} & \textbf{Γραμμές} & \textbf{Ευθύνη}\\
\midrule
\endfirsthead
\toprule
\textbf{Μονάδα} & \textbf{Γραμμές} & \textbf{Ευθύνη}\\
\midrule
\endhead
\multicolumn{3}{@{}l}{\emph{Έμπιστη βάση (\en{root})}}\\
\code{gate.py}          &  266 & κύκλος ζωής της καθολικής πύλης, επαναφορά\\
\code{gate\_server.py}  & 1409 & η προνομιούχος πύλη υπό \code{--privsep}\\
\code{protocol.py}      &  185 & το συμβόλαιο μηνυμάτων πύλης--αναλυτή\\
\code{privsep.py}       &  652 & \code{fork}, απόρριψη προνομίων, \en{subreaper}\\
\code{securefs.py}      &  160 & ασφαλείς λειτουργίες αρχείων (\code{O\_NOFOLLOW})\\
\code{safety.py}        &  384 & \en{watchdog}, αρχείο πανικού, εξαιρέσεις θυρών\\
\midrule
\multicolumn{3}{@{}l}{\emph{Ανάλυση (αναλυτής)}}\\
\code{daemon.py}        & 2011 & ο κεντρικός βρόχος γεγονότων και η ροή αποφάσεων\\
\code{sysfs.py}         &  869 & ανάγνωση του \en{sysfs}, καρφιτσωμένοι κατάλογοι\\
\code{descriptors.py}   &  334 & ανάλυση περιγραφέων \en{USB}\\
\code{descriptors\_\-safe.py} & 358 & αμυντικός περιηγητής αλυσίδας περιγραφέων\\
\code{usbclass.py}      &  288 & ταξινόμηση κλάσεων και λειτουργιών\\
\code{rules.py}         & 1057 & όλοι οι κανόνες και η διαμόρφωσή τους (\en{YAML})\\
\code{analyzers.py}     &  355 & ενορχήστρωση των σταδίων ανάλυσης\\
\code{quarantine.py}    &  369 & στάδιο 3: \code{EVIOCGRAB}, καταγραφή γεγονότων\\
\code{storage.py}       &  956 & στάδιο 4: \en{MBR}/\en{GPT}, συστήματα αρχείων\\
\code{mediawatch.py}    &  683 & παρακολούθηση καρτών μνήμης\\
\code{textsafe.py}      &  347 & καθαρισμός εχθρικού κειμένου\\
\code{ledger.py}        &  956 & βιβλίο ιστορικού, αποτυπώματα, μετάλλαξη\\
\code{trust.py}         &  588 & χώρος εμπιστοσύνης\\
\midrule
\multicolumn{3}{@{}l}{\emph{Διεπαφή χρήστη}}\\
\code{\_\_main\_\_.py}  &  830 & γραμμή εντολών\\
\code{agentlink.py}     &  885 & υποδοχή (\en{socket}) πράκτορα, έλεγχος \code{SO\_PEERCRED}\\
\code{agent.py}         &  771 & πράκτορας επιφάνειας εργασίας\\
\code{dialogs.py}       &  531 & \en{kdialog}, \en{zenity}, \en{tkinter}\\
\code{report.py}        &  273 & μορφοποίηση αναφοράς\\
\bottomrule
\end{longtable}

\section{Ο κύκλος ζωής της πύλης}
\label{sec:impl-gate}

Η κλάση \code{AuthorizationGate} υλοποιεί την πύλη ως διαχειριστή περιβάλλοντος
(\en{context manager}). Το Απόσπασμα~\ref{lst:gate-enter} δείχνει τον πυρήνα της
μεθόδου \code{\_\_enter\_\_}: πρώτα εγκαθίστανται οι χειριστές σημάτων και το
\code{atexit} — \emph{πριν} κλείσει οποιοσδήποτε διανομέας — ώστε να μην
υπάρχει στιγμή όπου η πύλη είναι κλειστή χωρίς διαδρομή επαναφοράς.

\begin{lstlisting}[language=Python,caption={Απλοποιημένη μορφή της μεθόδου εισόδου της πύλης (\code{gate.py}).},label=lst:gate-enter]
def __enter__(self) -> "AuthorizationGate":
    hubs = sysfs.list_root_hubs()
    if not hubs:
        raise RuntimeError("no USB root hubs found")
    self._armed = True
    self._owner_pid = os.getpid()     # only this process may reopen
    self._install_handlers()          # SIGTERM, SIGINT, SIGHUP
    atexit.register(self.restore)     # restore paths exist BEFORE closing
    for hub in hubs:
        current = sysfs.get_authorized_default(hub)
        if current is None:
            continue
        # A hub already at 0 is residue of a crashed run, not a choice:
        # restore 1, never 0, or the lockout becomes permanent.
        self._original[hub] = 1 if current == 0 else current
        sysfs.set_authorized_default(hub, 0)
    return self
\end{lstlisting}

Μια λεπτή περίπτωση αφορά τη διεργασία-παιδί του σταδίου~4. Η ανάγνωση του δίσκου
γίνεται σε διεργασία που δημιουργείται με \code{fork()}, η οποία κληρονομεί το
οπλισμένο αντικείμενο πύλης \emph{και} τους χειριστές σημάτων. Σε μια πρώιμη
έκδοση, όταν ο γονέας τερμάτιζε ένα «κολλημένο» παιδί με \code{SIGTERM}, ο
χειριστής του παιδιού έγραφε \code{authorized\_default=1} σε όλους τους
διανομείς, ενώ ο γονέας πίστευε ότι η πύλη ήταν κλειστή — ένα σφάλμα
\emph{αποτυχίας-ανοιχτά} (\en{fail-open}) το οποίο μια εχθρική συσκευή θα
μπορούσε να χρονομετρήσει. Η διόρθωση είναι το πεδίο \code{\_owner\_pid}: μόνο η
διεργασία που έκλεισε την πύλη μπορεί να την ανοίξει (βλ.~§\ref{sec:eval-qa}).

\section{Ανάλυση περιγραφέων}
\label{sec:impl-descriptors}

Το αρχείο \code{descriptors} του \en{sysfs} περιέχει τον περιγραφέα συσκευής
ακολουθούμενο από τις πλήρεις αλυσίδες διαμορφώσεων. Ο αναλυτής
(\code{descriptors\_safe.py}) διασχίζει την αλυσίδα αμυντικά: κάθε περιγραφέας
αρχίζει με \code{bLength} και \code{bDescriptorType}· μηδενικό ή υπερβολικό
μήκος, αλυσίδα συντομότερη από το \code{wTotalLength}, ή λιγότερες διαμορφώσεις
από όσες δηλώνει το \code{bNumConfigurations}, δεν προκαλούν εξαίρεση αλλά
\emph{ευρήματα} (\code{descriptor-chain-truncated}, \code{configurations-missing},
\code{descriptor-length-overstated}).

\section{Η μηχανή κανόνων}
\label{sec:impl-rules}

Κάθε κανόνας είναι ένα τμήμα κώδικα που, δεδομένου του περιβάλλοντος
(\code{Context}) της συσκευής, προσθέτει ένα εύρημα (\code{Finding}) με
αναγνωριστικό, σοβαρότητα, τίτλο και επεξήγηση. Η σοβαρότητα περνά από τη
διαμόρφωση του χρήστη (\code{cfg.severity}), ώστε ένα αρχείο \en{YAML}
(\code{--rules}) να μπορεί να απενεργοποιήσει ή να αναβαθμίσει/υποβαθμίσει κανόνες
των σταδίων 2–4. Οι κανόνες ιστορικού και ο \code{payload-captured} δεν
τροποποιούνται.

\begin{lstlisting}[language=Python,caption={Ο κανόνας της «υπογραφής \en{BadUSB}» (\code{probolos/rules.py}).},label=lst:rule-badusb]
# -- 1. The BadUSB signature ------------------------------------------
if has_keyboard and CLS_MASS_STORAGE in classes:
    add("storage-with-keyboard", Severity.CRITICAL,
        "Storage device that can also type",
        "This device presents itself as both a disk and a keyboard. The "
        "storage half provides the alibi while the keyboard half can enter "
        "commands on your behalf. Legitimate flash drives never need a "
        "keyboard interface. This is the classic BadUSB layout.")

# -- 2. Keystroke injection plus a network path ------------------------
if has_keyboard and (classes & {CLS_CDC, CLS_CDC_DATA, CLS_WIRELESS}):
    add("network-with-keyboard", Severity.CRITICAL, ...)
\end{lstlisting}

Το Απόσπασμα~\ref{lst:rule-timing} δείχνει την υλοποίηση του κριτηρίου
\eqref{eq:machine} του σταδίου~3.

\begin{lstlisting}[language=Python,caption={Διάκριση μηχανικής πληκτρολόγησης (\code{probolos/rules.py}).},label=lst:rule-timing]
_ACCIDENT_THRESHOLD = 5      # a lean on the keyboard, not a payload
_MACHINE_INTERVAL = 0.050    # 50 ms between keys: faster than sustained human
_MACHINE_REGULARITY = 0.20   # coefficient of variation below this is inhuman

intervals = obs.intervals()          # kernel timestamps, not Python's
mean_gap = _mean(intervals)
cv = _coefficient_of_variation(intervals)
machine_like = (
    len(keys) >= 3
    and mean_gap is not None and mean_gap < _MACHINE_INTERVAL
    and cv is not None and cv < _MACHINE_REGULARITY
)
if machine_like:
    add("machine-generated-keystrokes", Severity.CRITICAL, ...)
elif len(keys) >= _ACCIDENT_THRESHOLD:
    add("unprompted-typing", Severity.CRITICAL, ...)
else:
    add("unexpected-keystrokes", Severity.WARNING, ...)
\end{lstlisting}

\section{Η καραντίνα (στάδιο 3)}
\label{sec:impl-quarantine}

Η καραντίνα ζητά από την πύλη την ενεργοποίηση της συσκευής \emph{ανά
στιγμιότυπο}, και στη συνέχεια αναζητά τους κόμβους εισόδου της καθώς
εμφανίζονται, δεσμεύοντάς τους αμέσως με \code{EVIOCGRAB}. Νέοι κόμβοι
αναζητούνται καθ' όλη τη διάρκεια του παραθύρου. Μια αποτυχημένη δέσμευση
τερματίζει την παρατήρηση. Οι ενδιάμεσες μνήμες γεγονότων είναι φραγμένες, ώστε
μια συσκευή που «πλημμυρίζει» με γεγονότα να μην εξαντλεί τη μνήμη.

Χωρίς τη σημαία \code{--capture-payload}, ένα πάτημα πλήκτρου καταγράφεται
\emph{μόνο ως χρονοσφραγίδα}, όχι ως ταυτότητα πλήκτρου. Με τη σημαία, οι
πληκτρολογήσεις ανακατασκευάζονται σε απομαγνητοφώνηση τύπου
\en{DuckyScript} (διάταξη \en{US QWERTY}, έως περίπου 4096 χαρακτήρες), χρήσιμη
για ψηφιακή εγκληματολογία. Οι έλεγχοι εγγυώνται ότι (α) συσκευές που υπήρχαν
πριν από την εκκίνηση δεν επιθεωρούνται ποτέ και (β) μια εγκεκριμένη συσκευή δεν
καταγράφεται ποτέ — σημαντικό και για τη συμμόρφωση με τον \en{GDPR}.

\section{Επιθεώρηση αποθήκευσης (στάδιο 4)}
\label{sec:impl-storage}

Το Απόσπασμα~\ref{lst:mbr} δείχνει την ανάλυση του πίνακα \en{MBR}. Η ίδια
προσέγγιση εφαρμόζεται για την κεφαλίδα και τις εγγραφές \en{GPT}, ενώ οι
υπογραφές συστημάτων αρχείων αναγνωρίζονται από σταθερό σύνολο, χωρίς χρήση
του \en{libblkid}. Για τα \en{ISO~9660} και \en{UDF} απαιτείται και έλεγχος
δομής, όχι μόνο «μαγικός αριθμός».

\begin{lstlisting}[language=Python,caption={Ανάλυση των πρωτευόντων διαμερισμάτων \en{MBR} (\code{storage.py}).},label=lst:mbr]
def parse_mbr(data: bytes) -> List[Partition]:
    if len(data) < SECTOR:
        return []
    if struct.unpack_from("<H", data, 510)[0] != MBR_SIGNATURE:
        return []
    partitions = []
    for i in range(4):
        offset = 446 + i * 16
        status, _c1, _c2, _c3, ptype, _c4, _c5, _c6, start, sectors = \
            struct.unpack_from("<BBBBBBBBII", data, offset)
        if ptype == 0 or sectors == 0:
            continue        # empty slot
        partitions.append(Partition(index=i, bootable=bool(status & 0x80),
                                    type_byte=ptype, start_lba=start,
                                    sectors=sectors))
    return partitions
\end{lstlisting}

Η ανάγνωση εκτελείται σε διεργασία-παιδί με αυστηρό χρονικό όριο, ενώ ο
\en{watchdog} τίθεται σε παύση. Και τα δύο είναι απαραίτητα: το όριο μόνο του θα
μετρούσε την ανάγνωση ως «κόλλημα» του δαίμονα, ενώ η παύση μόνη της θα
μετέτρεπε ένα κόλλημα σε μόνιμο πάγωμα. Ο γονέας περιμένει την \emph{αναφορά}
του παιδιού και όχι την έξοδό του· η συσκευή απενεργοποιείται αμέσως μόλις
φτάσει η αναφορά.

\section{Διαχωρισμός προνομίων}
\label{sec:impl-privsep}

\subsection{Το πρωτόκολλο}

Το \code{protocol.py} ορίζει το πλήρες σύνολο αιτημάτων που μπορεί να εκτελέσει
η πύλη (Πίνακας~\ref{tab:protocol}). Κάθε μήνυμα είναι αντικείμενο \en{JSON} σε
ένα \en{datagram} \code{SEQPACKET}· το \code{SEQPACKET} διατηρεί τα όρια των
μηνυμάτων, ώστε να μη χρειάζεται πρόθεμα μήκους — μια ολόκληρη κατηγορία σφαλμάτων
ανάλυσης εξαλείφεται από την προνομιούχο διεργασία.

\begin{table}[htbp]
\centering
\caption{Τα αιτήματα του πρωτοκόλλου πύλης--αναλυτή. Ό,τι δεν περιλαμβάνεται
απορρίπτεται.}
\label{tab:protocol}
\begin{tabular}{@{}lp{10.2cm}@{}}
\toprule
\textbf{Αίτημα} & \textbf{Ενέργεια}\\
\midrule
\code{AUTHORIZE}   & θέτει \code{authorized=1/0} σε μία διαδρομή συσκευής ή διεπαφής\\
\code{SET\_DEFAULT} & θέτει \code{authorized\_default} σε έναν ριζικό διανομέα (επαναφορά μόνο στην καταγεγραμμένη τιμή)\\
\code{OPEN\_INPUT} & ανοίγει κόμβο \code{/dev/input/eventN} μόνο για ανάγνωση και επιστρέφει τον περιγραφέα\\
\code{OPEN\_BLOCK} & ανοίγει ολόκληρο δίσκο \code{/dev/sdX} μόνο για ανάγνωση και επιστρέφει τον περιγραφέα\\
\code{TRUST}       & προσθέτει εγγραφή εμπιστοσύνης, μόνο για συσκευή που μόλις ενεργοποίησε η πύλη και μόνο με το δικό της αποτύπωμα\\
\code{PING}        & έλεγχος ζωντάνιας\\
\bottomrule
\end{tabular}
\end{table}

\subsection{Η απόρριψη προνομίων}

Το Απόσπασμα~\ref{lst:drop} δείχνει τη συνάρτηση \code{drop\_privileges}. Η
σειρά (ομάδες $\rightarrow$ \en{gid} $\rightarrow$ \en{uid}) είναι η μόνη ορθή:
μετά το \code{setuid}, η διεργασία δεν θα είχε πλέον δικαίωμα να αλλάξει
ομάδες. Η επαλήθευση των συμπληρωματικών ομάδων προστέθηκε μετά από έλεγχο
κώδικα που εντόπισε ότι, σε ορισμένα περιβάλλοντα (\en{user namespaces}), μια
ομάδα μπορεί να επιβιώσει του \code{setgroups([])}.

\begin{lstlisting}[language=Python,caption={Επαληθευμένη απόρριψη προνομίων (\code{probolos/privsep.py}, χωρίς σχόλια).},label=lst:drop]
def drop_privileges(uid: int, gid: int) -> None:
    if os.getuid() != 0:
        return
    os.setgroups([])           # supplementary groups first
    os.setgid(gid)             # then group
    os.setuid(uid)             # then user, last of all
    if os.getuid() != uid or os.geteuid() != uid:
        raise PrivsepError("uid did not drop as expected")
    if os.getgid() != gid or os.getegid() != gid:
        raise PrivsepError("gid did not drop as expected")
    remaining = set(os.getgroups()) - {gid}
    if remaining:
        raise PrivsepError(f"supplementary groups survived: {sorted(remaining)}")
    try:
        os.setuid(0)
        raise PrivsepError("regained root after drop -- refusing to continue")
    except PermissionError:
        pass  # exactly what we want: root is unreachable now
\end{lstlisting}

\subsection{Κλειδί εμπιστοσύνης}

Το κλειδί εμπιστοσύνης (Απόσπασμα~\ref{lst:trustkey}) συνδυάζει ταυτότητα και
\en{SHA-256} των ακατέργαστων περιγραφέων. Μια συσκευή της οποίας οι περιγραφείς
δεν διαβάστηκαν δεν έχει κλειδί και, εκ κατασκευής, δεν μπορεί να γίνει έμπιστη.

\begin{lstlisting}[language=Python,caption={Υπολογισμός του κλειδιού εμπιστοσύνης (\code{probolos/trust.py}).},label=lst:trustkey]
def descriptor_hash(dev) -> Optional[str]:
    raw = getattr(dev, "raw_descriptors", None)
    return hashlib.sha256(raw).hexdigest() if raw else None

def identity_of(dev) -> str:
    serial = getattr(dev, "serial", None) or "-"
    return f"{dev.vendor_id}:{dev.product_id}:{serial}"

def key_for(dev) -> Optional[str]:
    digest = descriptor_hash(dev)
    if not digest:
        return None          # unreadable => untrustable by construction
    return f"{identity_of(dev)}#{digest}"
\end{lstlisting}

\section{Ασφαλής διαχείριση αρχείων κατάστασης}
\label{sec:impl-state}

Ο Πίνακας~\ref{tab:state-files} δείχνει τα αρχεία κατάστασης. Η εμπιστοσύνη και
το ιστορικό βρίσκονται σκόπιμα σε \emph{διαφορετικούς καταλόγους}: δικαίωμα
εγγραφής σε κατάλογο επιτρέπει την αντικατάσταση οποιουδήποτε αρχείου του,
ανεξάρτητα από τον ιδιοκτήτη του αρχείου. Μόνο ο \code{state/} παραδίδεται στον
αναλυτή.

\begin{table}[htbp]
\centering
\caption{Αρχεία κατάστασης, ιδιοκτησία και δικαιώματα.}
\label{tab:state-files}
\begin{tabular}{@{}p{5.6cm}p{3.6cm}p{4.2cm}@{}}
\toprule
\textbf{Διαδρομή} & \textbf{Περιεχόμενο} & \textbf{Ιδιοκτήτης / δικαιώματα}\\
\midrule
\pth{/var/lib/probolos/trusted.json} & εμπιστευμένες συσκευές & \en{root}· \code{0600} (\code{0644} υπό \code{--privsep})\\
\pth{/var/lib/probolos/state/ledger.json} & ιστορικό & αναλυτής· \code{0600}\\
\pth{/var/lib/probolos/instance.lock} & μία πύλη τη φορά & \en{root}\\
\pth{/run/probolos/agent.sock} & υποδοχή πράκτορα & \code{0660}, κατάλογος \code{2750}\\
\pth{/run/probolos.panic} & διακόπτης πανικού & πρέπει να το δημιουργήσει ο \en{root}\\
\bottomrule
\end{tabular}
\end{table}

Όλες οι εγγραφές γίνονται με \code{O\_NOFOLLOW | O\_CREAT | O\_EXCL} σε
προσωρινό αρχείο με απρόβλεπτο όνομα και στη συνέχεια ατομική μετονομασία
(\code{rename}), ώστε ένας συμβολικός σύνδεσμος να μην μπορεί να ανακατευθύνει
την εγγραφή. Κάθε αρχείο επικυρώνεται κατά τη φόρτωση· αρχεία άνω των 8\,\en{MiB},
ειδικά αρχεία ή μη έγκυρο \en{JSON} απορρίπτονται. Το ιστορικό διατηρείται κάτω
από 4\,\en{MiB}: όταν ξεπεραστεί, ξεχνιέται το ένα τέταρτο των εγγραφών,
ξεκινώντας από συσκευές για τις οποίες δεν απαντήθηκε ποτέ ερώτηση, ώστε μια
συσκευή που αλλάζει σειριακό αριθμό σε κάθε σύνδεση να μην μπορεί να «ξεπλύνει»
μια πραγματική απόρριψη.

\section{Ο πράκτορας επιφάνειας εργασίας}
\label{sec:impl-agent}

Ο πράκτορας (\code{agent.py}) εκτελείται στη συνεδρία του χρήστη και συνδέεται
στην υποδοχή \code{/run/probolos/agent.sock}. Ο αναλυτής ελέγχει το \en{uid} του
συνδεόμενου μέσω \code{SO\_PEERCRED}. Ο κατάλογος έχει δικαιώματα \code{2750}
και ιδιοκτήτη τον \en{root}· η υποδοχή δεσμεύεται από τον \en{root}
\emph{πριν} από την απόρριψη προνομίων και ο αναλυτής την κληρονομεί. Έτσι καμία
άλλη διεργασία με το ίδιο \en{uid} με τον αναλυτή δεν μπορεί να αντικαταστήσει
την υποδοχή και να «υποδυθεί» τον αναλυτή.

Οι διάλογοι υλοποιούνται με \en{kdialog} (\en{KDE}), \en{zenity} (\en{GNOME}) ή
\en{tkinter}. Μόνο οι \emph{δικές τους} απαντήσεις μετρούν: ένα κουμπί ή το
κλείσιμο του παραθύρου είναι απάντηση, ενώ ένας διάλογος που τερματίστηκε από
σήμα, κατέρρευσε ή έληξε είναι «καμία απάντηση». Η αντίστοιχη απεικόνιση
κωδικών εξόδου κάθε εργαλείου αποτέλεσε πηγή σφαλμάτων· χαρακτηριστικά, μια
έκδοση αντιστοίχιζε το «Όχι» του \en{tkinter} σε «Αποδοχή πάντα» (σφάλμα P1,
§\ref{sec:eval-qa}).

\section{Υπηρεσία \en{systemd}}
\label{sec:impl-systemd}

Το σενάριο \code{install.sh} εγκαθιστά τον κώδικα στο \code{/opt/probolos}
(ιδιοκτησίας \en{root}), δημιουργεί τον λογαριασμό \code{probolos} μέσω
\en{sysusers.d} και δύο μονάδες \en{systemd}: την πύλη ως υπηρεσία συστήματος
(\code{--privsep --agent}) με περιορισμούς απομόνωσης (π.χ.
\code{ProtectSystem=strict}, περιορισμένο σύνολο \en{capabilities}) και τον
πράκτορα ως υπηρεσία χρήστη. Όταν η επιφάνεια εργασίας δεν είναι διαθέσιμη
(οθόνη σύνδεσης), οι νέες συσκευές κρατούνται και ερωτώνται μόλις ο πράκτορας
επανασυνδεθεί.
